\section{Plethysm coefficients: the positivity and saturation indices}\label{sec:indices}

We keep the notation of Section~\ref{sec:plethysm}, with $n\ge4$, and put
$\zeta=e^{2\pi i/M}$, $\mathcal C_n(\rho)=\lim_{t\to\rho}(1-t^M)H_n(t)$ and
$A_n(\rho)=\mathcal C_n(\rho)/M$ for primitive $M$-th roots $\rho$. Since
$H_n$ has rational coefficients, $\mathcal C_n(\zeta^a)$ is the Galois
conjugate of $\mathcal C_n(\zeta)$ under $\zeta\mapsto\zeta^a$, and
$A_n(\zeta^{-a})=\overline{A_n(\zeta^a)}$. Define the $M$-periodic
\emph{dyadic wave}
\begin{equation}\label{eq:wave}
 W_n(r)=\sum_{\substack{1\le a<M\\a\ \mathrm{odd}}}A_n(\zeta^a)\zeta^{-ar}
 =\frac1M\Tr_{\Q(\zeta)/\Q}\bigl(\mathcal C_n(\zeta)\zeta^{-r}\bigr)\in\Q .
\end{equation}

\subsection{The wave and the indices}

By \cite[Theorem 1.6.1(b), p.~21]{GCT6}, $f_n$ is a quasi-polynomial, that is,
$H_n$ is a proper rational function with poles at roots of unity. So
\begin{equation}\label{eq:fourier}
 f_n(s)=\sum_\rho P_\rho(s)\rho^{-s}\qquad(s\ge0),
\end{equation}
the sum over the poles $\rho$ of $H_n$, with $P_\rho$ a polynomial of degree
one less than the pole order. At a simple pole $\rho$ of order $M$,
$1-t^M=(1-t/\rho)\prod_{\rho'\ne\rho}(1-t/\rho')$ over the other $M$-th roots
of unity $\rho'$, and $\prod_{\rho'\ne\rho}(1-\rho/\rho')=M$; so near $\rho$,
$H_n(t)=(\mathcal C_n(\rho)/M)/(1-t/\rho)+O(1)$, and $P_\rho=A_n(\rho)$ is
constant.

\begin{proposition}\label{prop:wave}
Let $L$ be the least common multiple of the orders of the poles of $H_n$, a
period of $f_n$, and write $f_n(s)=q_r(s)$ for $s\equiv r\pmod L$. Then
$v_2(L)=n-1$, and for every residue $r$ and every $X$,
\begin{equation}\label{eq:waveidentity}
 q_r(X)-q_{r+L/2}(X)=2W_n(r).
\end{equation}
Consequently $f_n(x)\ge2W_n(x)$ for every integer $x\ge p(f_n)$ and every
integer $x\ge s(f_n)+1$. In particular, if $2W_n(x)>f_n(x)$ for an integer
$x\ge0$, then $p(f_n)>x$ and $s(f_n)\ge x$.
\end{proposition}

\begin{proof}
By Theorem~\ref{thm:plethysmPH3}(2),(3), every pole order has two-part at most
$M$, and $M$ occurs; so $v_2(L)=n-1$. By \eqref{eq:fourier},
$q_r(X)=\sum_\rho P_\rho(X)\rho^{-r}$. If the order of $\rho$ has two-part at
most $M/2$, it divides $L/2$ and $\rho^{-L/2}=1$; if it has two-part $M$, then
$\rho=\zeta^a$ with $a$ odd (Theorem~\ref{thm:plethysmPH3}(3)) and
$\rho^{-L/2}=-1$, as $L/M$ is odd. Subtracting the two expansions gives
\eqref{eq:waveidentity}. If $p=p(f_n)$, every constituent polynomial of
$n\mapsto f_n(n+p)$, namely $X\mapsto q_r(X+p)$, has nonnegative coefficients,
so $q_r(y)\ge0$ for all real $y\ge p$ and all $r$. If $s=s(f_n)$, every
constituent $X\mapsto q_r(X+s)$ is either identically zero or positive at every
integer $X\ge1$, so $q_r(y)\ge0$ for every integer $y\ge s+1$. At such $x$, take
$r\equiv x\pmod L$ in \eqref{eq:waveidentity}: $f_n(x)-2W_n(x)=
q_{r+L/2}(x)\ge0$.
\end{proof}

\subsection{An exact product for the residues}

Write $w_i=(-2)^i$ for $0\le i\le n-2$ and $w_{n-1}=0$, and form the multisets
\begin{equation}\label{eq:Sn}
\begin{gathered}
 \mathcal E_n=\{\{0,1,k\}:2\le k<n\}\cup\{\{0,2,3\}\},\\
 \mathcal S_n=\{w_i:0\le i\le n-2\}\ \sqcup\ \{w_i+w_j+w_k:0\le i<j<k<n,\
 \{i,j,k\}\notin\mathcal E_n\},
\end{gathered}
\end{equation}
with repetitions retained. Then $|\mathcal S_n|=(n-1)+\binom n3-(n-2)-1=
\binom n3$, and no element of $\mathcal S_n$ is divisible by $M$: a singleton
$w_i$ ($i\le n-2$) has valuation $i<n-1$, and in a sum of distinct $w$'s the
smallest valuation among the nonzero terms occurs once.

\begin{proposition}\label{prop:product}
For every odd $a$,
\begin{equation}\label{eq:product}
 \mathcal C_n(\zeta^a)=\prod_{b\in\mathcal S_n}\bigl(1-\zeta^{a\nu b}\bigr)^{-1}.
\end{equation}
\end{proposition}

\begin{proof}
By Galois conjugation it suffices to take $a=1$. Put $\eta=\zeta^\nu$ and
$\alpha_i=\eta^{w_i}$, and let $D=\diag(\alpha_0,\dots,\alpha_{n-1})$, which
fixes $F_n$ (its chain monomials have exponents $2w_i+w_{i+1}\equiv0$ and
$3w_{n-1}=0$). On coordinate functions of $\Sym^3V^*$, that is on
$\Sym(\Sym^3V)$, $D$ acts through the character of $\Sym^3V$, and on invariants
of coefficient degree $d_0s$ by $(\det D)^{\gamma s}$. Here
$(\det D)^\gamma=\zeta^{3\sum_iw_i}$ and $3\sum_iw_i=1-(-2)^{n-1}\equiv1
\pmod M$; so $D$ acts on $T^{(n)}$ as the grading element $\zeta$. By the
tangent space computation in Proposition~\ref{prop:chaincubic}(4), the
cotangent space of $X$ at $p_n$ has the character $\Sym^3V-(V\otimes V^*-1)$
as a representation of $\langle D\rangle$, and the conormal space of $Y$ is
obtained by removing the trivial character of the fixed direction:
\begin{equation}\label{eq:conormal}
 N^*=\Sym^3V-V\otimes V^* .
\end{equation}
With $h_j$, $e_j$ the complete and elementary symmetric functions in the
$\alpha_i$, the identity $\prod_i(1+t\alpha_i)=\prod_i(1-t^2\alpha_i^2)/
\prod_i(1-t\alpha_i)$ gives $h_3=e_3+h_1\sum_i\alpha_i^2$. The multisets
$\{\alpha_i^2\}$ and $\{\alpha_i^{-1}\}$ differ only in that $\alpha_0^{-1}$ is
replaced by $1$, because $\alpha_i^2=\alpha_{i+1}^{-1}$ for $i<n-1$ and
$\alpha_{n-1}=\alpha_{n-2}^2=1$. So \eqref{eq:conormal} is
$e_3+h_1-\alpha_0^{-1}h_1$. The trivial weight $\alpha_{n-1}$ in $h_1$ cancels
the term $i=0$ of $\alpha_0^{-1}h_1$. For $2\le k<n$ the subtracted weight
$w_k-w_0$ equals the weight of the triple $\{0,1,k\}$, as $w_0+w_1=-1$; the
last subtracted weight $w_1-w_0=-3$ equals that of $\{0,2,3\}$ ($1+4-8$ for
$n\ge5$; $1+4+0\equiv-3\pmod8$ for $n=4$). These triples are distinct, and what
remains is the multiset $\mathcal S_n$. Equation \eqref{eq:plethysmresidue}
gives \eqref{eq:product}.
\end{proof}

\subsection{A dominant pair}

Put $r_\nu=3\nu\in\{3,9\}$, let $1\le a_*<M$ with $r_\nu a_*\equiv1\pmod M$
(so $a_*$ is odd and $a_*\not\equiv-a_*$), and
\begin{equation}\label{eq:ZE}
 Z_n=A_n(\zeta^{a_*}),\qquad
 E_n=\sum_{\substack{a\bmod M\ \mathrm{odd}\\a\not\equiv\pm a_*}}|A_n(\zeta^a)| .
\end{equation}

\begin{proposition}\label{prop:amplitude}
For every $n\ge4$,
\begin{equation}\label{eq:amplitude}
 \log_2|Z_n|\ge\binom n4-3\binom{n-1}3-(n-1)(n-2)-\binom n3-(n-1)
 \ge\frac{n^4}{24}-2n^3 .
\end{equation}
If $n\ge6$ and $E_n\le|Z_n|/4$, then for every integer $X$ there is an integer
$x$ with $X\le x\le X+8$ and $W_n(x)\ge|Z_n|/4$.
\end{proposition}

\begin{proof}
By \eqref{eq:product}, $Z_n=M^{-1}\prod_{b\in\mathcal S_n}(1-\theta^b)^{-1}$ with
$\theta=\zeta^{\nu a_*}$, so that $\theta^3=\zeta$. Every factor has absolute
value at most two. Consider the triples $i<j<k\le n-2$ of $\mathcal S_n$, those
not using the zero weight. Their weight $w=w_i+w_j+w_k$ is divisible by three
(each $w_i\equiv1\pmod3$) and $|w|<2^{k+1}$, so
$|1-\theta^w|=|1-\zeta^{w/3}|\le2\pi|w|/(3M)\le8\cdot2^k/M$, and
$-\log_2|1-\theta^w|\ge(n-1-k)-3$. Summing $n-1-k$ over all triples
$i<j<k\le n-2$ counts the choices of a fourth index in $\{k+1,\dots,n-1\}$ and
gives $\binom n4$; at most $n-2$ of these triples belong to $\mathcal E_n$, each
with a term at most $n-1$. The remaining at most $\binom n3$ factors contribute
at least $-1$ each, and $M^{-1}$ contributes $-(n-1)$. This proves the first
inequality of \eqref{eq:amplitude}. For the second use $\binom n4\ge
n^4/24-n^3/4$, $3\binom{n-1}3\le n^3/2$, $\binom n3\le n^3/6$ and
$(n-1)(n-2)+(n-1)\le n^2\le n^3/4$: the total cubic loss is at most
$\tfrac76n^3$.

Write $r_\nu a_*=1+kM$; then $1\le k<r_\nu$ and $\gcd(k,r_\nu)=1$, since
$kM\equiv-1\pmod{r_\nu}$. Thus $a_*/M=k/r_\nu+1/(r_\nu M)$. For fixed $X$, the
$r_\nu$ numbers $Z_n\zeta^{-a_*X}e^{-2\pi ikj/r_\nu}$, $0\le j<r_\nu$, are
equally spaced on the circle of radius $|Z_n|$, so one has real part at least
$\cos(\pi/r_\nu)|Z_n|\ge|Z_n|/2$. Replacing the phase $e^{-2\pi ikj/r_\nu}$ by
$\zeta^{-a_*j}$ changes it by a factor $e^{-2\pi ij/(r_\nu M)}$ and the real
part by at most $(2\pi/M)|Z_n|\le|Z_n|/4$ for $n\ge6$. So
$2\operatorname{Re}(Z_n\zeta^{-a_*(X+j)})\ge|Z_n|/2$ for this $j\le8$. The terms
$a=\pm a_*$ of \eqref{eq:wave} are this conjugate pair, and the others have
total absolute value at most $E_n\le|Z_n|/4$.
\end{proof}

\begin{corollary}\label{cor:conditional}
If $E_n\le|Z_n|/4$ for all sufficiently large $n$, then for every real
$0<c<1/4$ and all sufficiently large $n$, $p(f_n)>2^{cn}$ and $s(f_n)\ge2^{cn}$.
\end{corollary}

\begin{proof}
Take $X=\lceil2^{cn}\rceil$ and $x\le X+8$ as in
Proposition~\ref{prop:amplitude}. The value $f_n(x)$ is at most the dimension
of $\Sym^{d_0x}(\Sym^3V)$, so
$\log_2\max\{1,f_n(x)\}\le(m-1)\log_2(m+nx)=\frac c6n^4+O_c(n^3\log n)$, which
is eventually smaller than $\log_2(2W_n(x))\ge n^4/24-2n^3-1$, as $c/6<1/24$.
Proposition~\ref{prop:wave} gives $p(f_n)>x\ge2^{cn}$ and $s(f_n)\ge x$.
\end{proof}

\subsection{Reduction to an increment}

For $n\ge4$ write $M_n=2^{n-1}$, $b_n=3^{-1}\bmod M_n$, and for odd $a$
\begin{equation}\label{eq:Pn}
 P_n(a)=\prod_{b\in\mathcal S_n}|1-e^{2\pi iab/M_n}|^{-1},\qquad
 Q_n=\sum_{\substack{a\bmod M_n\ \mathrm{odd}\\a\ne\pm b_n}}\frac{P_n(a)}{P_n(b_n)} .
\end{equation}
By \eqref{eq:product}, $|A_n(\zeta^a)|=P_n(\nu a)/M_n$; multiplication by the
odd number $\nu$ permutes the odd residues, and $\nu a_*\equiv b_n$. So
$Q_n=E_n/|Z_n|$ and $P_n(b_n)=M_n|Z_n|$. For $n\ge5$ form the multiset of
$\binom{n-1}2$ residues
\begin{equation}\label{eq:Tn}
 \mathcal T_n=\{1\}\sqcup\{1+w_j+w_k:2\le j<k<n,\ (j,k)\ne(2,3)\}
 \sqcup\{2+w_k:3\le k<n\}\sqcup\{6\},
\end{equation}
and $J_n(a)=\prod_{b\in\mathcal T_n}|1-e^{2\pi iab/M_n}|^{-1}$.

\begin{proposition}\label{prop:recursion}
For $n\ge5$ and odd $a$, $P_n(a)=J_n(a)\,P_{n-1}(-a\bmod M_{n-1})$. The selected
sibling ratio satisfies
\begin{equation}\label{eq:sibling}
 \delta_n:=\frac{P_n(b_n+M_{n-1})}{P_n(b_n)}\le2^{-n^3/6+4n^2-107n/6+19},
 \qquad\delta_n\le2^{-n^3/12}\ (n\ge48).
\end{equation}
If for some $0<\rho<1$ and all $n\ge n_0$ and all odd $a$ with
$a\bmod M_{n-1}\notin\{\pm b_{n-1}\}$
\begin{equation}\label{eq:contraction}
 J_n(a)+J_n(a+M_{n-1})\le\rho J_n(b_n),
\end{equation}
then $Q_n\le\rho Q_{n-1}+2\delta_n$ for $n\ge n_0$, and $Q_n\to0$.
\end{proposition}

\begin{proof}
The weights $w_1,\dots,w_{n-1}$ are $-2$ times the weights of the chain in
$n-1$ variables, in the same order. All triples with smallest index at least one
lie in $\mathcal S_n$; the scaled $\mathcal S_{n-1}$ omits among them exactly
$\{1,2,k\}$ ($3\le k<n$) and $\{1,3,4\}$; and $\mathcal S_n$ has the extra
singleton $w_0=1$ and the triples $\{0,j,k\}$, $2\le j<k<n$, other than
$\{0,2,3\}$. Their weights are those of \eqref{eq:Tn}: $w_1+w_2+w_k=2+w_k$, and
$\{1,3,4\}$ has weight $-2-8+16=6$ ($n\ge6$) or $-2-8+0\equiv6\pmod{16}$ ($n=5$).
The count is $1+(\binom{n-2}2-1)+(n-3)+1=\binom{n-1}2$. Substituting
$e^{2\pi i a(-2b')/M_n}=e^{2\pi i(-a)b'/M_{n-1}}$ gives the recursion.

The residues $b_n$ and $b_n+M_n/2$ have the same smaller factor, and every even
weight gives the same factor at both; so $\delta_n=\prod_{b\in\mathcal O_n}
|\tan(\pi b_nb/M_n)|$ over the odd weights $\mathcal O_n=\{1\}\sqcup\{1+w_j+w_k\}$.
The bulk ones, $2\le j<k\le n-2$, $(j,k)\ne(2,3)$, number
$B=\binom{n-3}2-1$, are divisible by three with $|b|<2^{k+1}\le M_n$; as
$3b_n\equiv1$, $|\tan(\pi b_nb/M_n)|=|\tan(\pi b/(3M_n))|\le2\pi|b|/(3M_n)\le
8\cdot2^k/M_n$, using $|\sin u|\le|u|$ and $\cos u\ge1/2$ on $[-\pi/3,\pi/3]$.
The sum of $n-1-k$ over them is $\binom{n-2}3-(n-4)$. The other $n-2$ odd
factors are at most $M_n/2$, an odd residue being at distance at least one from
$M_n/2$, so the cosine is at least $\sin(\pi/M_n)\ge2/M_n$. Hence
$\log_2\delta_n\le3B-\binom{n-2}3+(n-4)+(n-2)^2=-n^3/6+4n^2-107n/6+19$, which
is at most $-n^3/12$ for $n\ge48$.

Group the odd residues modulo $M_n$ by their reduction modulo $M_{n-1}$. The
two selected parents $\pm b_{n-1}$ have the selected children $\pm b_n$ and the
children $\pm b_n+M_{n-1}$, of total normalized mass $2\delta_n$ (conjugation).
Every other parent has two children whose normalized mass is at most $\rho$
times its own in $Q_{n-1}$, by the recursion and \eqref{eq:contraction}, since
$P_n(b_n)=J_n(b_n)P_{n-1}(b_{n-1})$. So $Q_n\le\rho Q_{n-1}+2\delta_n$, and as
$\delta_n\to0$, $Q_n\to0$.
\end{proof}

\subsection{The increment estimates}

\begin{lemma}\label{lem:sparse}
Let $1\le r\le3$, integers $1\le q_1<\dots<q_r$, signs $\epsilon_i$,
$\kappa\in\Z$, $J\ge1$ odd, $R=\lceil\log_2(16J)\rceil$, and
$x=\kappa/3+(J/3)\sum_i\epsilon_i2^{-q_i}$. Then $\|x\|_{\mathbb R/\Z}\ge
1/(6\cdot2^{q_1+(r-1)R})$. In particular, if $q_1\le R$, then
$2^{-3R-1}\le|1-e^{2\pi ix}|\le2$.
\end{lemma}

\begin{proof}
Keep the initial block of terms ending before the first gap $q_{i+1}-q_i>R$
(all terms if there is none), ending at $q_t$. Its numerator over $3\cdot
2^{q_t}$ is odd (the last term is odd, the earlier ones and $\kappa2^{q_t}$
even), so it is at distance at least $1/(3\cdot2^{q_t})$ from $\Z$. The tail
is less than $(2J/3)2^{-q_{t+1}}\le J2^{-R}/(3\cdot2^{q_t})\le
1/(48\cdot2^{q_t})$. So $\|x\|\ge1/(6\cdot2^{q_t})$, and $q_t\le q_1+(r-1)R$.
Finally $|1-e^{2\pi ix}|\ge4\|x\|$.
\end{proof}

Let $l=n-1$, $M=2^l$, and for odd $J$ let $b_J\equiv J\cdot3^{-1}\pmod M$; then
$\eta_J=\zeta^{b_J}$ satisfies $\eta_J^3=\zeta^J$ and
\begin{equation}\label{eq:cuberoot}
 \eta_J=\omega_Je^{2\pi iJ/(3M)},\qquad\omega_J^3=1,\qquad
 \omega_J=1\iff3\mid J\quad(J<M),
\end{equation}
since $3b_J=J+k_JM$ with $k_J\in\{0,1,2\}$, $\omega_J=e^{2\pi ik_J/3}$, and
$k_J=0$ only if $3\mid J$. For $|y|\le1/24$,
$|\log_2(\sin(\pi y)/(\pi y))|\le8y^2$, by $1-u^2/6\le\sin u/u\le1$.

\begin{proposition}[near increments]\label{prop:near}
For odd $1\le j<M$, with $a_j\equiv j\cdot3^{-1}\pmod M$ and
$R=\lceil\log_2(16j)\rceil$, and $n\ge5$,
\begin{equation}\label{eq:near}
 \Bigl|\log_2\frac{J_n(a_j)}{J_n(b_n)}+\binom{n-1}2\log_2j
 -\mathbf1_{3\mid j}\binom n2\Bigr|\le100nR^2 .
\end{equation}
For $\delta=10^{-5}$ and odd $3\le j\le2^{\delta n}$ this gives
$\log_2(J_n(a_j)/J_n(b_n))\le-n^2\log_2j/12$.
\end{proposition}

\begin{proof}
The weights of $\mathcal T_n$ are: the \emph{bulk} weights $1+w_u+w_v$
($2\le u<v\le l-1$, $(u,v)\ne(2,3)$), the weights $2+w_k$ ($3\le k\le l-1$) and
$6$, all divisible by three, $B=(n-1)(n-4)/2$ in number; and the
\emph{exceptional} weights $1$, $2$ and $1+w_k$ ($2\le k\le l-1$), $n-1$ in
number. Give each weight its largest active dyadic index $k$ and put $q=l-k$
(the weights $1$ and $2$ get $q=l$ and $q=l-1$). The exceptional weights have
each $q=1,\dots,l$ exactly once, so $B+(n-1)=\binom{n-1}2$ and the exceptional
$q$ sum to $\binom n2$. A weight divisible by three has absolute value between
$2^{k-2}$ and $2^{k+1}$ (for $1+w_u+w_v$ use $2^v-2^{v-1}-1\ge2^{v-2}$, $v\ge3$);
for each $q$ there are at most $2n$ of them. An exceptional weight has absolute
value between $2^{k-1}$ and $3\cdot2^{k-1}$. By \eqref{eq:cuberoot}, the target
root is $\omega_je^{2\pi ij/(3M)}$. A factor is \emph{low} if $q>R$ and
\emph{high} otherwise. Lemma~\ref{lem:sparse} applies to every factor (each
weight is a signed sum of at most three powers of two): a high target factor is
at least $2^{-3R-1}$, and at least $2^{-2R-1}$ for an exceptional one; all
factors are at most two.

For a low factor divisible by three there is no cube-root shift: its target
phase is $j$ times its reference phase and is at most $(2j/3)2^{-q}<1/24$ in
absolute value; by the sine estimate the logarithm of the factor ratio differs
from $\log_2j$ by at most $16j^24^{-q}$; summing over at most $2n$ weights for
each $q>R$ gives less than $n$, as $j^24^{-R}\le1/256$. For a high factor
divisible by three, the reference phase lies between $2^{-q}/12$ and
$2\cdot2^{-q}/3<1/2$ in absolute value, so the reference factor is at least
$2^{-q}/3$, and the logarithm of the ratio is at most $5R$ in absolute value;
there are at most $2nR$ of them, and restoring their $\log_2j$ main terms costs
at most $2nR^2$. So the factors divisible by three contribute $-B\log_2j$ to the
logarithm of the reciprocal product, with error at most $20nR^2$.

Every exceptional reference factor lies between $1/3$ and two: its cube-root
shift is nontrivial and its phase perturbation is at most $1/4$. If $3\nmid j$,
every low exceptional target factor has the same kind of shift, perturbed by
less than $1/24$, so its logarithmic ratio is at most three; the at most $R$
high ones are at most $3R$ each; the total is at most $10nR^2$. If $3\mid j$,
the target shift vanishes; a low exceptional target factor lies between
$(2/3)j2^{-q}$ and $4j2^{-q}$, so its reciprocal ratio has logarithm $q-\log_2j$
up to four. Restoring the $q$ of the at most $R$ high factors costs at most
$R^2$, their actual ratios at most $3R^2$, discarding the $-\log_2j$ terms at
most $nR$, and the low errors at most $4n$; the exceptional contribution is
$\binom n2$ up to $10nR^2$. Altogether
$\log_2(J_n(a_j)/J_n(b_n))=-B\log_2j+\mathbf1_{3\mid j}\binom n2+\mathcal E$ with
$|\mathcal E|\le30nR^2$, and replacing $B$ by $\binom{n-1}2$ costs
$(n-1)\log_2j\le nR$.

For the consequence, put $t=\log_2j>3/2$, so $R\le t+5$. With
$\binom{n-1}2\ge n^2/2-2n$ and $\binom n2\le n^2/2$, \eqref{eq:near} gives at
most $-(t-1)n^2/2+n(2t+100(t+5)^2)\le-tn^2/6+2502nt^2$, using $t-1\ge t/3$ and
$t+5\le5t$; and $t\le n/10^5$ with $2502/10^5<1/12$.
\end{proof}

\begin{proposition}[far increments]\label{prop:far}
Let $n\ge5$, $l=n-1$, $M=2^l$, $a$ odd, $j\equiv3a\pmod M$ with $1\le j<M$, and
$J=\min\{j,M-j\}$. Then
\begin{equation}\label{eq:far}
 \log_2\frac{J_n(a)}{J_n(b_n)}\le7n^2-\frac{(\log_2J-8)_+^3}{24}.
\end{equation}
In particular, for $\delta>0$, $J\ge2^{\delta n}$ and
$n\ge\max\{5,16/\delta,2688/\delta^3\}$, the ratio is at most
$2^{-\delta^3n^3/384}$.
\end{proposition}

\begin{proof}
Consider the \emph{active pair} factors $1-\exp(2\pi ia(1+w_j+w_k)/M)$,
$2\le j<k<l$; the increment contains all but $(j,k)=(2,3)$. Every phase is
nonzero modulo one, its numerator being odd. Put $\theta=a/M$,
$y_{jk}=\theta(1+w_j+w_k)$ modulo one, and for $0\le u\le l-1$ let $r=l-u$ and
$N_u$ be the number of active pairs with $\|y_{jk}\|<2^{-u}$.

The map sending a subset of $\{0,\dots,l-1\}$ to the sum of its $w_i$ modulo
$2^l$ is injective (a difference of two subset sums has a unique least index
with coefficient $\pm1$, of valuation less than $l$). Call an index early if
$j<r$ and late otherwise; there are $e=(r-2)_+$ early indices. There are at
most $\binom e2$ early--early pairs. For a fixed early $j$, the mixed numerators
$a(1+w_j+w_k)$ are congruent modulo $2^r$ and distinct modulo $M$; the condition
$\|y_{jk}\|<2^{-u}$ puts their centred representatives in $(-2^r,2^r)$, which
holds at most three of one class; so there are at most $3e$ good mixed pairs.
The late--late numerators are all congruent to $a$ modulo $2^r$ and distinct,
so at most three are good. Hence
\begin{equation}\label{eq:Nu}
 N_u\le\binom r2+3r+3 .
\end{equation}

Let $\tau=\|3\theta\|=J/M$ and suppose $\tau>3\cdot2^{-u}$. Partition the early
indices into $h=\lfloor e/2\rfloor$ consecutive pairs $(2,3),(4,5),\dots$. For
two of them, $(j,j+1)$ and $(k,k+1)$, the identity
$y_{jk}+y_{j+1,k}+y_{j,k+1}=3\theta\pmod1$ (from $w_{j+1}=-2w_j$) shows that
the three phases cannot all be within $2^{-u}$ of an integer; so one of the four
pairs of the block is bad. The $\binom h2$ blocks are disjoint and do not
contain $(2,3)$. So under this condition
$N_u\le\binom r2+3r+3-\binom{\lfloor(r-2)_+/2\rfloor}2$.

For a nonzero phase $y$ with denominator dividing $M$, $\|y\|\ge1/M$ and
$\log_2|1-e^{2\pi iy}|^{-1}\le-\log_2\|y\|\le\sum_{u=0}^{l-1}
\mathbf1_{\{\|y\|<2^{-u}\}}$. The other $2n-4$ increment factors have
logarithmic reciprocals at most $l$ each, less than $2n^2$ in total. Let
$K=(\lfloor\log_2J\rfloor-2)_+$; for $1\le r\le K$, $J>3\cdot2^r$, so the
condition holds at $u=l-r$. Summing,
\[
 \log_2J_n(a)\le\binom n3+\tfrac32l(l+1)+3l+2n^2
 -\sum_{r=1}^K\binom{\lfloor(r-2)_+/2\rfloor}2 ,
\]
where the first four terms are at most $n^3/6+4n^2$ for $n\ge5$, and the sum is
at least $(K-5)_+^3/24\ge(\log_2J-8)_+^3/24$, since for $r\ge6$ the summand is
at least $(r-3)(r-5)/8\ge(r-5)^2/8$.

For the selected increment, the $B'=\binom{n-3}2-1$ bulk active pairs
(divisible by three) have factor at most $8\cdot2^k/M$ at $a=b_n$, as in the
proof of Proposition~\ref{prop:recursion}, and all other factors are at most
two. So
\[
 \log_2J_n(b_n)\ge\binom{n-2}3-(n-4)-2B'-\binom{n-1}2
 =\frac{n^3}6-3n^2+\frac{71n}6-11\ge\frac{n^3}6-3n^2 .
\]
Subtracting gives \eqref{eq:far}. If $\log_2J\ge\delta n\ge16$, then
$(\log_2J-8)^3/24\ge\delta^3n^3/192$, and $7n^2\le\delta^3n^3/384$ for
$n\ge2688/\delta^3$.
\end{proof}

\begin{theorem}\label{thm:indices}
$E_n/|Z_n|\to0$. Consequently, for every real $0<c<1/4$ and all sufficiently
large $n$,
\[
 p(f_n)>2^{cn},\qquad s(f_n)\ge2^{cn},
\]
and $p(f_n)>2^{n/8}$, $s(f_n)\ge2^{n/8}$ for every $n\ge2\cdot10^{80}$.
\end{theorem}

\begin{proof}
For an odd residue $a\ne\pm b_n$, let $j$ be the least positive residue of $3a$
up to sign, $1<j\le M/2$ (and $j=1$ exactly for $a=\pm b_n$); as
$J_n(a)=J_n(-a)$, Proposition~\ref{prop:near} covers $j\le2^{\delta n}$ with
ratio at most $2^{-n^2\log_23/12}$, and Proposition~\ref{prop:far} the rest with
ratio at most $2^{-\delta^3n^3/384}$, for $n$ large. So the two nonselected
children of any parent other than $\pm b_{n-1}$ have normalized increments of
sum at most $1/2$ eventually: \eqref{eq:contraction} holds with $\rho=1/2$, and
Proposition~\ref{prop:recursion} gives $Q_n=E_n/|Z_n|\to0$.
Corollary~\ref{cor:conditional} gives the first assertion.

For the explicit bound, let $N=10^{20}$. For $n\ge N$ both ratios are at most
$2^{-n^2/8}$ ($\log_23>3/2$; and $\delta^3n^3/384\ge n^2/8$ for
$n\ge48/\delta^3$), so \eqref{eq:contraction} holds with $\rho=1/2$ for all
$n\ge N$. Every factor of $P_n(a)$ is at least $4/M_n$, so
$P_n(a)\le2^{(n-3)\binom n3}$, while $P_n(b_n)=M_n|Z_n|\ge2^{-2n^3}$ by
\eqref{eq:amplitude}; hence $Q_n\le2^{n-1+(n-3)\binom n3+2n^3}\le2^{n^4}$.
Iterating from $N$, with $\delta_k\le2^{-k^3/12}$,
\[
 Q_n\le2^{-(n-N+1)}Q_{N-1}+2\sum_{k=N}^n2^{-(n-k)}\delta_k
 \le2^{-n+N-1+N^4}+4\cdot2^{-N^3/12}\le\frac14\qquad(n\ge2N^4).
\]
For such $n$ take $X=\lceil2^{n/8}\rceil$ in Proposition~\ref{prop:amplitude}:
some $x\in[X,X+8]$ has $2W_n(x)\ge2^{n^4/24-2n^3-1}$. As $\log_2n\le n/16$ for
$n\ge128$, $2^{n/8}\ge n^2$; also $m\le n^3/5$ for $n\ge16$ and $x\le2^{n/8+1}$,
so $m+nx\le4n2^{n/8}$ and
\[
 \log_2\max\{1,f_n(x)\}\le(m-1)\log_2(m+nx)\le\frac{n^3}5\Bigl(\frac n8+
 \log_2n+2\Bigr)\le\frac{3n^4}{80}+\frac{2n^3}5<\frac{n^4}{24}-2n^3-1
\]
for $n\ge1024$. Proposition~\ref{prop:wave} gives $p(f_n)>x\ge2^{n/8}$ and
$s(f_n)\ge x$.
\end{proof}

\begin{corollary}\label{cor:PH2SH}
Hypotheses 1.6.5(a) (SH) and 1.6.6 (PH2) of \cite{GCT6} fail for the plethysm
coefficients of the ray $\mu=(3)$, $\lambda=(d_0)$, $\pi=((3/\nu)^n)$: the
indices exceed every polynomial in $h=m^2+n+1=\Theta(n^6)$. Hypotheses
3.3.2(a) and 3.3.3 fail for the plethysm problem with $\rank(x)=2n+m+1$.
\end{corollary}

